\section{\ref{sec:nora}장. \nt{NORA}}

\nt{NORA} 시스템의 구조, 두 작동 모드, 동공과의 연관(\nt{NORA}만을 통한 것은 아니다), 신호 대 배경 비의 증가, 행동에서의 뒤집힌 U는 직접 측정되었다. 곱셈 이득 $g$는 모델이다. 이득이 선택 모드를 전환하고 역온도로 작동한다는 것은 이 장의 유도 결과다. 이득 조정의 이점은 이 책의 모델로 계산한 것이다. “\nt{NORA}는 탐색의 노브다”와 “\nt{NORA}는 인터럽트다”는 가설이다.

{\footnotesize
\setlength{\tabcolsep}{3pt}\renewcommand{\arraystretch}{1.15}
\begin{longtable}{>{\raggedright}p{0.5cm}>{\raggedright}p{4.0cm}>{\raggedright}p{5.6cm}>{\raggedright}p{2.3cm}>{\raggedright\arraybackslash}p{1.7cm}}
\toprule
No. & 명제 & 실험과 측정한 것 & 출처 & 상태 \\
\midrule\endfirsthead
\toprule
No. & 명제 & 실험과 측정한 것 & 출처 & 상태 \\
\midrule\endhead
1 & 사람의 \nt{NORA} 뉴런은 수만 개이며, 거의 모두 한 무리에 있다 & 사람 뇌줄기의 연속 절편, 티로신 수산화효소 염색: 청반에 $53\,900$개, 이웃한 하위핵에 $6\,260$개의 세포. & \cite{Baker1989LC} & 측정됨 \\
2 & 출력은 거의 뇌 전체로 퍼지지만, 무리의 부분들이 서로 다른 영역을 어느 정도 나누어 맡는다 & 생쥐에서의 바이러스 유전 표지: 청반의 \nt{NORA} 뉴런은 뇌의 넓은 영역에서 입력을 받고, 뇌와 척수의 넓은 영역으로 출력을 보낸다. 쥐: 편도체로 가는 무리는 공포 학습을 강화하고, 전전두피질로 가는 별도의 무리는 공포 소거를 강화한다. & \cite{Schwarz2015LC, Uematsu2017LC, Poe2020} & 측정됨 \\
3 & 배경 모드: 초당 1회 미만에서 몇 회 사이의 고른 스파이크, 발화율은 각성 수준에 따라 변한다 & 쥐, 수면--각성 주기 동안의 청반 뉴런 기록: 발화율은 깨어 있을 때 가장 높고, 서파 수면에서 낮으며, 렘수면에서는 거의 0이다. 발화율 변화가 상태 전환보다 앞선다. & \cite{AstonJonesBloom1981} & 측정됨 \\
4 & 돌발 모드: 과제에 중요한 사건에 대략 결정 순간에 나오는 짧은 버스트 & 원숭이, 드문 표적 찾기 과제: 모든 청반 뉴런이 표적에만 버스트로 응답한다. 표적 뒤 $\sim90\ms$, 손 반응 $\sim200\ms$ 전이다. 수행이 나쁠 때 버스트가 약하다. 양자택일 과제에서 버스트는 자극보다 반응에 더 정확히 맞물리며, 정답과 오답 모두에서 나타난다. & \cite{AstonJones1994vigilance, Clayton2004LC} & 측정됨 \\
5 & 동공 지름은 \nt{NORA}의 부정확한 측정점이다. \nt{NORA}뿐 아니라 \nt{ACET}와 다른 영역도 따른다 & 원숭이: 청반의 스파이크는 세포 하나의 개별 스파이크까지 동공 확장을 예측한다. 비슷하지만 더 늦은 연관이 둔덕과 띠피질에도 있다. 생쥐: 동공 변동은 피질로 가는 노르아드레날린성 출력과 콜린성 출력의 활동을 모두 따른다. & \cite{Joshi2016, Reimer2016} & 측정됨 \\
6 & 노르아드레날린은 배경 활동을 유발 활동보다 더 강하게 억제한다. 곱셈 이득~\eqref{eq:gain}은 이 효과를 전달하는 모델이다 & 깨어 있는 원숭이, 청각 피질, 노르아드레날린 이온영동: 뉴런의 배경 활동이 같은 종의 소리에 대한 응답보다 더 강하게 억제되어 신호 대 잡음비가 커진다. 곱셈 시그모이드~\eqref{eq:gain}는 네트워크 모델에서 가져왔다. 기울기가 문자 그대로 곱해진다는 것은 측정되지 않았다. & \cite{Foote1975NE, ServanSchreiber1990} & 측정됨; $g$는 모델 \\
7 & 이득은 증폭기 뒤의 잡음에 대해서만 $d'$를 높인다~\eqref{eq:gainsnr}(명제~\ref{prop:gainnoise}) & 이 장의 유도. 네트워크 모델에서 이득은 신호 검출을 개선한다. 증폭기 앞과 뒤의 잡음을 분리해 이 차이를 검증한 실험은 없다. & 이 장의 유도; \cite{ServanSchreiber1990} & 이론 \\
8 & 경계 $g\beta=1$은 선택 네트워크를 “숙고 중”에서 “결정함”으로 전환한다~\eqref{eq:wtagain}(명제~\ref{prop:gainwta}, 그림~\ref{fig:pitchfork}) & 서로 억제하는 두 무리의 선형 분석; 이 장의 유도. 뇌에서 이 문턱값을 직접 측정한 것은 없다. & 이 장의 유도 & 이론 \\
9 & 이득은 선택의 역온도다~\eqref{eq:softmax} & 증폭기 뒤의 잡음이 로지스틱이면 선택 확률은 $T=\sigma/g$인 softmax다. 이 장의 유도. & 이 장의 유도 & 이론 \\
10 & \nt{NORA}는 얼마나 탐색할지 정한다. 높은 배경 활동은 작은 유효 $g$(탐색), 결정 순간의 돌발 버스트는 큰 $g$(결정) & 사람: 무작위 선택 직전의 기저 동공은 알려진 최선을 고르기 직전보다 넓다. 동공이 넓은 사람일수록 탐색 성향이 높다. 쥐: 무작위 선택 모드로의 전환은 앞띠피질로 가는 \nt{NORA} 입력이 조절한다. 버스트가 유효 $g$를 높인다는 것은 직접 측정되지 않았다. & \cite{JepmaNieuwenhuis2011, Tervo2014stochastic, AstonJones2005} & 가설 \\
11 & 보상은 $g$에 대해 뒤집힌 U로 의존한다. 빠르게 변하는 세계에서 최선의 $g$가 더 작다(명제~\ref{prop:explore}, 그림~\ref{fig:invertedU}) & \texttt{models/nora\_explore.py}, $4000$번 시도의 실행 $60$회. $1000$번에 한 번 변화: 최선 $g=16$, 비율 $0.976$; $20$번에 한 번: $g=8$, 비율 $0.886$; $g=0.5$에서 $0.77$. 재계산 결과가 본문과 일치했다. & \texttt{models/\allowbreak nora\_\allowbreak explore.py} & 모델 \\
12 & 행동에서의 뒤집힌 U: 배경 활동이 적당하고 버스트가 뚜렷할 때 가장 잘한다 & 원숭이, 4행과 같은 과제: 수행이 좋은 시기에는 배경 발화율이 적당하고 표적에 대한 버스트가 뚜렷하다. 배경 발화율이 높으면 버스트가 없고 오경보가 많다. 배경 활동과 유발 활동의 변화는 수행의 변동을 따라간다. & \cite{Usher1999LC, AstonJones2005} & 측정됨 \\
13 & \nt{NORA}는 예상치 못한 불확실성을, \nt{ACET}는 예상된 불확실성을 알린다 & 두 종류의 불확실성을 다루는 베이즈 추론 이론. 인용한 연구에는 이 신호들을 신경전달물질별로 분리한 직접 실험이 없다. & \cite{YuDayan2005} & 가설 \\
14 & 급격한 변화 뒤 사람은 더 빨리 학습하고 동공이 커진다 & 사람, 갑작스러운 변화가 있는 예측 과제: 동공의 짧은 변화는 변화 확률의 추정을 따르며, 기저 동공과 함께 새 데이터가 추정을 얼마나 바꾸는지(학습률)를 예측한다. 빛이나 과제와 무관한 동공 변화도 이 학습률을 바꾼다. 여기서 동공이 바로 \nt{NORA}를 보여 준다는 것은 5행의 간접 자료에 기댄 추론이다. & \cite{Nassar2012} & 측정됨 \\
15 & 돌발 버스트는 인터럽트로 작동한다. 네트워크가 상태를 초기화한다 & \nt{NORA} 버스트가 네트워크 상태를 초기화한 직접 실험은 없다. 측정된 것은 중요한 사건에 대한 버스트뿐이다(4행). & \cite{BouretSara2005} & 가설 \\
16 & 뜻밖의 일에 따라 $g$나 학습률을 조정해도 최선의 고정 설정보다 거의 낫지 않다 & \texttt{models/nora\_explore.py}, 시기당 $1000$번 시도인 세계($1000$번에 한 번, $20$번에 한 번 변화). 최선의 고정 $g=8$: $0.925$; $g$ 조정: 최대 $0.927$; 학습률 조정: 최대 $0.926$; 고정 학습률 $0.4$: $0.928$. 재계산 결과가 본문과 일치했다. & \texttt{models/\allowbreak nora\_\allowbreak explore.py} & 모델 \\
\bottomrule
\end{longtable}}
